
\begin{definition}[Kleene star]
    For a set $\cS \subseteq \bbN^n$,
    the \emph{Kleene star} of $\cS$ is defined as
    \begin{equation*}
        \kstar{\cS}\ :=\ \setc{\sum_{s \in \cS'} s}{\text{$\cS'$ is a finite multisubset of $\cS$}}.
    \end{equation*}
    Note that $\bfzero \in \kstar{\cS}$ for any $\cS$.
\end{definition}

If $\cS$ is semilinear then its Kleene star is semilinear: this is closely-related to Parikh's theorem which gives a correspondence between semi-linear sets and regular languages: see \cite{parikh,kleenestar}.
We will be interested in bounding the complexity of the Kleene star of a semilinear set using numerical data associated to the set.
Note that for an integer linear system $\cQ(\bfx): \bfA \bfx = \bfc$,
if $\bfc = \bfzero$, then $\kstar{\sol{\cQ(\bfx)}} = \sol{\cQ(\bfx)}$.
The following result gives a bound when $\bfc \neq \bfzero$.

\begin{lemma}\label{lem:key}
    For every integer linear system $\cQ(\bfx): \bfA \bfx = \bfc$ with $\bfc \neq \bfzero$,
    where $\bfA \in \bbZ^{m \times n}$ and $\bfc \in \bbZ^{m}$,
    there exists $\bfAtilde \in \bbZ^{(n+t) \times (n+k)}$ with $\norm{\bfAtilde} = 1$
    such that
    $\kstar{\sol{\cQ(\bfx)}} = \sol{\cQtilde(\bfx)}$,
    where $\cQtilde(\bfx)$ is the  Boolean combination of integer linear systems
    \begin{equation*}    
        \exists \bfy_1, \bfy_2\ 
        (\bfx\ =\ \projp{n}{\bfy_1} + \bfy_2)\ \land\ 
        (\bfAtilde\bfy_1\ =\ \bfzero)\ \land\ 
        (\bfA\bfy_2\ =\ \bfzero)\ \land\ 
        ((\bfy_1=\bfzero) \to (\bfy_2=\bfzero)),
    \end{equation*}
    $k = n\cdot(2D+1)^{m} + 1$,
    $t \le 2nk$, and
    $D = n \cdot \norm{\bfA} \cdot \left((n+1)\cdot \norm{\bfA} + \norm{\bfc} + 1\right)^{m}$.
    Moreover, $\bfAtilde$ can be computed in time $2^\bigO{\log n + \log K + m^2}$,
    where $K := \max\left(\norm{\bfA}, \norm{\bfc}\right)$.
\end{lemma}

We defer the proof of this to the appendix.

